\documentclass[11pt,a4paper]{article}
\usepackage[T1]{fontenc}
\usepackage{lmodern,amsmath,amssymb,amsthm,microtype}
\usepackage[margin=25mm]{geometry}
\usepackage[colorlinks=true,urlcolor=blue,linkcolor=blue]{hyperref}
\newtheorem{proposition}{Proposition}
\newcommand{\dd}{\mathop{}\!\mathrm d}
\title{Combined vertices in the pure-recoil limit}
\author{Alan Day \and Chengheng\thanks{AI research assistant (OpenAI).}}
\date{Working research note --- 30 October 2026}
\begin{document}\maketitle
\begin{abstract}
We isolate the combined vertex \(V/\sqrt C\) in the internal integral
representation of the continued tetrahedron graph. An exact cancellation,
followed by two inverse-\(b\) shifts, gives its normalized pure-recoil
limit as an ordinary Gamma quotient. The doubled vertex has a simple
hyperbolic form. Away from mass-triangle degeneration, four such vertices
reduce to a mass polynomial, uniformly on a specified transverse box.
This does not identify the physical four-point normalization or complete
the internal integral. In particular this vertex must not be multiplied
onto a full-\(6j\) Gram asymptotic without redoing the prefactor bookkeeping.
All results are author self-reviewed drafts, without a priority claim.
\end{abstract}

\section{Which factors are being computed}
Set the de Sitter radius to one. In CLP~\cite{CLP}, (3.17) and
(3.43)--(3.47), write \(Q=b+b^{-1}\), \(p_a=iP_a\), and
\(x_\sigma=Q/2+\sum_{a=1}^3\sigma_a p_a\). Let
\(\mathcal E=\{\sigma:\sigma_1\sigma_2\sigma_3=+1\}\) and
\(\mathcal O=-\mathcal E\). Then
\begin{align}
 C&=K_b\,\frac{\prod_{\sigma\in\mathcal E\cup\mathcal O}
                   \Gamma_b(x_\sigma)}
                 {\prod_a\Gamma_b(Q+2p_a)\Gamma_b(Q-2p_a)},&
 K_b&=\frac{\Gamma_b(2Q)}{\sqrt2\,\Gamma_b(Q)^3},\\
 V^2&=\prod_{\sigma\in\mathcal E}S_b(x_\sigma),&
 S_b(x)&=\frac{\Gamma_b(x)}{\Gamma_b(Q-x)}.
\end{align}
Consequently, as a meromorphic identity,
\begin{equation}
 U_b:=\frac{V^2}{C}
 =K_b^{-1}\frac{\prod_a\Gamma_b(Q+2p_a)\Gamma_b(Q-2p_a)}
                   {\prod_{\sigma\in\mathcal O}\Gamma_b(x_\sigma)^2}.
 \label{eq:exact}
\end{equation}
Square-root branches belong to the original continued graph.
On the real Euclidean germ the doubled combined vertex is
\(T_b=|V/\sqrt C|^2=|U_b|\).

The whole graph can thus be written as
\begin{equation}
 Z_{\rm tet}=\left(\prod_{\text{four vertices}}T_b\right)
              |\mathcal J_b|^2, \qquad
 \{6j\}_b=\left(\prod_{\text{four vertices}}V_b\right)\mathcal J_b.
 \label{eq:book}
\end{equation}
Here \(\mathcal J_b\) is precisely the internal double-sine integral
in CLP (3.46), with its specified contour. This rearrangement does not
supply physical insertion legs or external energy cycles.

\section{A normalized pure-recoil vertex}
Use \(b^2=i/\kappa\) and the pure-recoil parametrization
\[
 p_a=-\frac1{2b}-\frac{ibm_a}{2},\quad m_a>0,\qquad
 L=b^{-2}=-i\kappa,\quad S=m_1+m_2+m_3 .
\]
Define
\[
 z_0=\frac{1+iS}{2},\qquad
 z_a=\frac{1-i(S-2m_a)}2\quad(a=1,2,3).
\]
The four odd arguments in \eqref{eq:exact} are \(2/b+bz_0,bz_1,bz_2,bz_3\).
The six leg arguments are \(b(1-im_a),2/b+b(1+im_a)\).
The exact shift relation gives
\begin{equation}
 \Gamma_b(2/b+bz)=
 \frac{2\pi\,b^{1-L-2z}}{\Gamma(z)\Gamma(L+z)}\Gamma_b(bz).
 \label{eq:shift}
\end{equation}
All powers of \(b\) generated in \eqref{eq:exact} have mass-independent
net exponent. Up to a factor depending only on \(b\), one obtains
\begin{equation}
 U_b(\mathbf m)=
 E_b(\mathbf m)
 \frac{\Gamma(z_0)^2\Gamma(L+z_0)^2}
      {\prod_a\Gamma(1+im_a)\Gamma(L+1+im_a)},\qquad
 E_b=\frac{\prod_{a,\pm}\Gamma_b(b(1\pm im_a))}
                 {\prod_{j=0}^3\Gamma_b(bz_j)^2}.
 \label{eq:shifted}
\end{equation}

\begin{proposition}
For fixed positive real masses \(\mathbf m\) and a fixed positive
reference triple \(\mathbf r\), with the branches inherited from CLP,
\begin{equation}
 \lim_{\kappa\to\infty}\frac{U_b(\mathbf m)}{U_b(\mathbf r)}
 =\frac{F(\mathbf m)}{F(\mathbf r)},\qquad
 F(\mathbf m)=
 \frac{\prod_a\Gamma(1-im_a)}
      {\prod_a\Gamma(\frac12-\frac i2(S-2m_a))^2}.
 \label{eq:F}
\end{equation}
In particular \(T_b(\mathbf m)/T_b(\mathbf r)\to T(\mathbf m)/T(\mathbf r)\),
where
\begin{equation}
 T(\mathbf m)=
 \frac1{\pi^{3/2}}\sqrt{\frac{m_1m_2m_3}
                   {\prod_a\sinh(\pi m_a)}}
 \prod_a\cosh\!\left(\frac{\pi(S-2m_a)}2\right).
 \label{eq:T}
\end{equation}
Only a mass-independent normalization has been suppressed.
\end{proposition}

\begin{proof}
The usual small-\(b\) reduction to Gamma functions is familiar from
the Schwarzian limit~\cite{MTV}. To control this particular balanced
quotient on \(b^2=i/\kappa\), a direct argument avoids assuming a
uniform continuation of that real-\(b\) formula.

Put \(f_\kappa(z)=\log\Gamma_b(bz)-\log\Gamma(z)\).
The third logarithmic derivative of the canonical Barnes product is
absolutely convergent and gives, on the pole-free domain in use,
\begin{equation}
 f_\kappa'''(z)=-2\sum_{q=1}^{\infty}\sum_{n=0}^{\infty}
                         (z+n-iq\kappa)^{-3}. \label{eq:lattice}
\end{equation}
This also follows by differentiating the Barnes double-zeta definition
three times; quadratic normalization terms disappear. For
\(\operatorname{Re}z\ge1/2\), \(|\operatorname{Im}z|\le R\), and
\(\kappa\ge2R\), comparison of the \(n\)-sum with an integral gives
\begin{equation}
 |f_\kappa'''(z)|\le
 K_\kappa:=\frac{16\zeta(3)}{\kappa^3}
                +\frac{8\zeta(2)}{\kappa^2}. \label{eq:bound}
\end{equation}
Indeed, with \(A=q\kappa/2\), the sum of
\(((n+1/2)^2+A^2)^{-3/2}\) is at most \(A^{-3}+A^{-2}\).

The signed arguments in \(E_b\), with coefficients \(+1\) for each
\(1\pm im_a\) and \(-2\) for each \(z_j\), satisfy
\begin{equation}
 \sum c= -2,\qquad \sum c z=2,\qquad \sum c z^2=4.
 \label{eq:moments}
\end{equation}
All three are independent of the masses. Taylor-expand \(f_\kappa\)
about \(z=1\) along line segments in the domain above. The quadratic
polynomial cancels between a triple and its reference. The remaining
logarithmic error is at most
\begin{equation}
 \frac{K_\kappa}{6}
 \left(\sum_{\mathbf m}|c|\,|z-1|^3+
       \sum_{\mathbf r}|c|\,|z-1|^3\right). \label{eq:remainder}
\end{equation}
Thus \(E_b(\mathbf m)/E_b(\mathbf r)\) has the corresponding ordinary
Gamma quotient as its limit. This establishes the needed normalized
reduction on this ray without choosing an absolute Barnes normalization.

Finally put \(B_L(\mathbf m)=\Gamma(L+z_0)^2/
\prod_a\Gamma(L+1+im_a)\). The fixed-shift Gamma ratio
asymptotics~\cite{DLMF} give
\[
 B_L(\mathbf m)\sim\Gamma(L)^{-1}L^{-2},
\]
since \(2z_0-\sum_a(1+im_a)=-2\). Its ratio to the reference tends to
one. Substitution into \eqref{eq:shifted} proves \eqref{eq:F}.
The identities
\(|\Gamma(1-im)|^2=\pi m/\sinh(\pi m)\) and
\(|\Gamma(1/2-ix/2)|^2=\pi/\cosh(\pi x/2)\) give \eqref{eq:T}.
\end{proof}

Equation \eqref{eq:remainder} controls the balanced Barnes part only.
The fixed-shift estimate for \(B_L\) has not here been replaced by a
fully quantified joint growing-mass bound. We therefore use the
sequential limit: first \(\kappa\to\infty\) at fixed masses, then the
large-mass regime below. No finite-\(G\) error estimate follows.

\section{A local transverse box, with an explicit vertex remainder}
For real positive masses satisfying strict triangle inequalities, set
\(\ell_a=S-2m_a>0\). Equation \eqref{eq:T} factors exactly as
\begin{equation}
 T(\mathbf m)=\frac{\sqrt{m_1m_2m_3}}{2\sqrt2\,\pi^{3/2}}
 \mathcal E(\mathbf m),\qquad
 \mathcal E=
 \frac{\prod_a(1+e^{-\pi\ell_a})}
      {\prod_a(1-e^{-2\pi m_a})^{1/2}}. \label{eq:E}
\end{equation}
The exponentials cancel because \(\sum_a\ell_a=S\).
Continue this expression from that real germ. In a domain with
\(\operatorname{Re}\ell_a\ge d>0\) and
\(\operatorname{Re}m_a\ge a>0\), the square root and logarithm fixed
there obey
\begin{equation}
 |\log\mathcal E|\le
 \frac{3e^{-\pi d}}{1-e^{-\pi d}}
 +\frac{3e^{-2\pi a}}{2(1-e^{-2\pi a})}. \label{eq:Ebound}
\end{equation}
This is an analytic continuation, not a new pointwise modulus at
complex masses.

For equal probes \(\nu\), write the cyclic observer masses
\[
 m_1=M+\omega/2+u,\quad m_3=M+\omega/2-u,\qquad
 m_2=M-\omega/2+v,\quad m_4=M-\omega/2-v.
\]
The four vertex triples are \((m_i,m_{i+1},\nu)\), cyclically.
Every graph edge occurs twice. Their product, normalized at
\(\omega=u=v=0\), is therefore
\begin{align}
 \mathcal R_V&=P\,\exp\{\Delta\log\mathcal E_{\rm four}\},\\
 P&=\frac{[(M+\omega/2)^2-u^2][(M-\omega/2)^2-v^2]}{M^4}.
 \label{eq:poly}
\end{align}
The logarithmic remainder is bounded by the sum of eight bounds
\eqref{eq:Ebound}, four at each configuration.

Take \(\omega=iA_0y\), \(A_0=\nu^2/(2M)\), \(|y|\le Y\), and
real \(|u|,|v|\le L_{\rm box}\sqrt M\). If
\[
 \nu/M\to0,\qquad L_{\rm box}\sqrt M=o(\nu),\qquad \nu\to\infty,
\]
then all the real parts of the triangle gaps are at least
\(\nu-2L_{\rm box}\sqrt M\), except for the much larger gaps
\(m_i+m_{i+1}-\nu\). The exponentials are uniformly small.
Moreover, with \(a=\omega/(2M),r=u/M,s=v/M\),
\[
 P-1=-2a^2+a^4-(r^2+s^2)(1+a^2)+2a(r^2-s^2)+r^2s^2 .
\]
It follows uniformly on this box that
\begin{equation}
 \mathcal R_V=1+
 O_Y\!\left(\frac{\nu^4}{M^4}+\frac{L_{\rm box}^2}{M}\right)
 +O\!\left(e^{-\pi(\nu-2L_{\rm box}\sqrt M)}\right). \label{eq:flat}
\end{equation}
For example the earlier hierarchy \(\nu^{4/3}\ll M\ll\nu^2\)
allows such boxes, with a slowly growing \(L_{\rm box}\).
This controls the limiting vertex on a local box; it does not bound
the full integration tail or prove the action's Gaussian approximation.

\section{Known observer factors and the remaining physical gap}
CLP (3.10) gives the observer factor
\(Z_{\rm obs}(\alpha)=32\sin^2(\alpha/2)\sinh^2(\kappa\alpha/2)\).
For \(\alpha_i=2\pi(1-m_i/\kappa)\) and \(\sum_i m_i=4M\), it implies
\begin{equation}
 \lim_{\kappa\to\infty}
 \frac{\prod_{i=1}^4Z_{\rm obs}(\alpha_i)}
      {Z_{\rm obs}(2\pi(1-M/\kappa))^4}=P^2. \label{eq:obs}
\end{equation}
The common entropy and Boltzmann exponent cancel at fixed total mass.
This is a check of a known building block, not a prescription that
the physical four-point function contains precisely this product.

The two-point insertion factor in CLP (3.37) is only a leading
fixed-probe prescription. Its displayed \((\alpha_1\alpha_2)^{-3/2}\)
dependence becomes constant in the sequential pure-recoil limit,
but its unspecified correction and the four-point insertion rule
are not determined by this observation. We do not impose them.

There is also a bookkeeping correction to guard against. The
previous Gram/transverse calculation used a \emph{full-\(6j\)}
geometric one-loop hypothesis. The \(V\)'s in \eqref{eq:book} are
already part of that \(6j\). Hence \(\mathcal R_V\) is not a new factor
to multiply into the previous Gram answer. It tells us that in the
\emph{internal-integral representation} the combined vertices are
locally flat after the specified limits. The determinant and order-one
factor of \(\mathcal J_b\) must be computed in that same representation
before the two organizations can be compared.

The next step is therefore the pure-recoil degeneration of the full
internal double-sine integral, retaining \(u,v\), followed by its
contour and saddle analysis. The physical external-energy cycle,
normalization, possible additional saddles and finite-\(\kappa\)
uniformity remain independent requirements. No full OTOC, microscopic
dictionary, or de Sitter dual has been established.

\section{Reproducibility}
The script \texttt{code/recoil\_vertex.py} records 42 checks at 65
decimal digits: three exact signed-moment identities, four
Gamma/hyperbolic identities and their remainder bounds, six
large-Gamma shift ratios, four lattice-derivative bounds, nine
complex transverse-box factorizations and bounds, and three observer
ratios. The six shift ratios and three observer ratios are repeated
at 95 digits. For the lattice derivative the \(n\)-sum is a polygamma
function and the \(q\)-sum is truncated at 40 with the explicit
absolute tail bound from \eqref{eq:bound}. No direct numerical
evaluation of the entire double-Gamma quotient or \(6j\) is claimed.
Finite samples do not certify the uniform statements or a physical
finite-\(G\) amplitude.

\begin{thebibliography}{9}
\bibitem{CLP} J. Chandra, D. Liska and B. Post,
\emph{3D de Sitter from complex Virasoro TQFT: Observers and knots},
\href{https://arxiv.org/html/2610.30247v1}{arXiv:2610.30247v1},
(3.10), (3.17), (3.35)--(3.47) and Appendix A.
\bibitem{MTV} T. G. Mertens, G. J. Turiaci and H. L. Verlinde,
\emph{Solving the Schwarzian via the conformal bootstrap},
\href{https://doi.org/10.1007/JHEP08(2017)136}{JHEP 08 (2017) 136},
Appendix B.1. Online reading only; no new local archive.
\bibitem{DLMF} NIST Digital Library of Mathematical Functions,
\href{https://dlmf.nist.gov/5.11}{Section 5.11},
fixed-shift Gamma asymptotics.
\end{thebibliography}
\end{document}
